% ==========================================================
% 同行、不同组与最低人数限制
% 难度：★★★ 难题
% ==========================================================

\newcommand{\qProbPermSpecGroupStem}{%
八名同学前往武当山、黄山、庐山三个不同景点，每人只选择一个景点，每个景点至少有两名同学。八人中有四名男生、四名女生，并满足：男生甲与女生 $A$ 不去同一景点；女生 $B$ 与女生 $C$ 去同一景点。求游玩行程的方法数。
}

\newcommand{\qProbPermSpecGroupAnswer}{%
$564$
}

\newcommand{\qProbPermSpecGroupSolution}{%
记男生甲为 $M$，女生 $A$ 为 $F$，女生 $B,C$ 必须同行，其余四人称为普通同学。先安排 $B,C,M,F$，再分配四名普通同学。

\textbf{情形一：$B,C$ 所在景点与 $M,F$ 所在景点均不同}
此时三个景点基础人数为 $(2,1,1)$。
特殊对象的安排方式：选择 $B,C$ 所在景点有 $3$ 种；$M,F$ 分别去另外两个景点有 $2$ 种。特殊对象共 $3\times 2=6$ 种。
四名普通同学原本有 $3^4$ 种，但两个当前只有一人的景点必须至少再得一人。利用容斥原理：$3^4-2\times 2^4+1 = 50$ 种。
该情形共有 $6\times 50=300$ 种。

\textbf{情形二：$B,C$ 与 $M,F$ 中的一人同去一个景点}
此时三个景点的基础人数为 $(3,1,0)$。
特殊对象的安排方式：选择 $B,C$ 所在景点 $3$ 种；选择 $M,F$ 谁与 $B,C$ 同行 $2$ 种；另一人去剩余两景点之一 $2$ 种。共 $3\times 2\times 2=12$ 种。
分配四名普通同学：设当前空景点得到 $r$ 名普通同学，为至少两人，需 $r\geqslant 2$。
当 $r=2$ 时：选两人去空景点 $\binom{4}{2}=6$ 种。剩余两人去三人或一人景点，但一人景点需至少再得一人，有 $2^2-1=3$ 种。共 $6\times 3=18$ 种。
当 $r=3$ 时：选三人去空景点 $\binom{4}{3}=4$ 种。剩余一人必须去一人景点，有 $1$ 种。共 $4\times 1=4$ 种。
当 $r=4$ 时：一人景点无法达到两人要求，不合法。
所以普通同学共 $18+4=22$ 种。
该情形共有 $12\times 22=264$ 种。

总方法数为 $300+264=564$ 种。
}

\newcommand{\qProbPermSpecGroupTeachingNotes}{%
\textbf{【避坑与边界说明】}
\begin{itemize}
    \item 三个景点互不相同，最终位置有标号。
    \item $B,C$ 是两个不同的人，只是必须选择同一景点，不能把二人误当成一个人计算容量。
    \item 男生甲与女生 $A$ 只要求不同景点，并没有规定具体去哪个景点。
    \item 安排好特殊对象后，必须根据各景点当前人数补足“至少两人”。
    \item 容斥时，减去某景点无人补入后，要把两个景点同时无人补入的情况加回来。
\end{itemize}
}
